% ---- score2 standard preamble (from: python3 score2_tools.py preamble) ----
\documentclass[10pt]{extarticle}
\usepackage[a4paper,margin=1.5cm]{geometry}
\usepackage[T1]{fontenc}
\usepackage{amsmath,amssymb}
\usepackage{graphicx}
\usepackage[export]{adjustbox}
\usepackage{multicol}
\usepackage{enumitem}
\usepackage{booktabs,array,tabularx}
\setlength{\columnsep}{0.7cm}
\setlength{\parindent}{0pt}
\setlength{\parskip}{3pt}
\pagestyle{empty}
\makeatletter
% figure registry (figdefs.tex is generated by the extract step)
\newcommand{\qfigdef}[4]{\@namedef{qf@file@#1}{#2}\@namedef{qf@w@#1}{#3}\@namedef{qf@wide@#1}{#4}}
\newcommand{\qf@check}[1]{\@ifundefined{qf@file@#1}{\PackageError{score2}{Unknown figure id #1}{}}{}}
\newcommand{\qf@img}[4][]{\includegraphics[width=#2\dimexpr\@nameuse{qf@w@#4}pt\relax,max width=#3,max height=0.45\textheight,#1]{\@nameuse{qf@file@#4}}}
% Full-width questions: finalize wraps every question that contains a wide figure in \qfullbegin...\qfullend,
% so the WHOLE question (stem, figures, options) is set across the page in reading order, then the two columns
% resume. The question's head - its number, stem and first figure (\qfullsplit marks the end of the head) -
% is kept on one page, so a number is never left at a page foot with its figure overleaf; the rest may break
% across pages normally, so no big gaps are left. Without a \qfullsplit, the whole question is kept together
% if it is under 45% of a page. Never write these yourself.
\newif\ifqf@full
\newif\ifqf@boxed
\newsavebox{\qf@box}
\newcommand{\qf@open}{\global\qf@fulltrue\par\global\qf@boxedtrue
  \setbox\qf@box\vbox\bgroup\hsize\textwidth\linewidth\textwidth}
\newcommand{\qf@emit}[1]{\par\egroup\global\qf@boxedfalse
  \ifdim\dimexpr\ht\qf@box+\dp\qf@box\relax<#1\textheight\noindent\usebox{\qf@box}\par\else\unvbox\qf@box\fi}
\newcommand{\qfullsplit}{\ifqf@boxed\qf@emit{0.85}\fi}
\newcommand{\qf@close}{\ifqf@boxed\qf@emit{0.45}\fi\global\qf@fullfalse}
\newcommand{\qfullbegin}{\end{multicols}\qf@open}
\newcommand{\qfullend}{\qf@close\begin{multicols}{2}}
\newcommand{\qfullstart}{\qf@open}
\newcommand{\qfulldone}{\qf@close}
% \qfig[scale]{id}: block figure, centred, directly where it is written (size set by the tool).
% Fallback if a wide figure is not inside a full-width question: break out of the columns in place - no floats.
\newcommand{\qfig}[2][1]{\qf@check{#2}%
  \ifqf@full\qf@blk{#1}{#2}\else
  \ifnum\@nameuse{qf@wide@#2}=1 \expandafter\@firstoftwo\else\expandafter\@secondoftwo\fi
  {\end{multicols}\begin{center}\qf@img{#1}{\textwidth}{#2}\end{center}\begin{multicols}{2}}%
  {\qf@blk{#1}{#2}}\fi}
\newcommand{\qf@blk}[2]{\par\smallskip{\centering\qf@img{#1}{\linewidth}{#2}\par}\smallskip}
% \qfigi[scale]{id}: inline figure (inside option lists / tables / minipages); never breaks columns.
\newcommand{\qfigi}[2][1]{\qf@check{#2}\qf@img[valign=c]{#1}{\linewidth}{#2}}
\makeatother
% structure macros - lint reads these to check completeness and figure placement
\newcommand{\papersec}[1]{\par\bigskip{\centering\Large\bfseries PAPER-#1\par}\medskip}
\newcommand{\partsec}[2]{\par\medskip\noindent{\large\bfseries PART-#1: #2}\par\smallskip}
\newcommand{\sect}[1]{\par\smallskip\noindent\textbf{#1}\par}
\newcommand{\qn}[1]{\par\smallskip\noindent\textbf{#1.}~}
\newcommand{\paraq}[2]{\par\smallskip\noindent\textbf{Paragraph for Questions #1 and #2}\par}
\newenvironment{qoptions}{\begin{itemize}[leftmargin=1.8em,itemsep=1pt,topsep=2pt,label={}]}{\end{itemize}}
\newcommand{\opt}[2]{\item[(#1)] #2}
% figure registry (generated by score2_tools.py finalize)
\qfigdef{p1_phy_s1_q03_1}{TOAS - Test 06 (17 Apr 2022)-img/p1_phy_s1_q03_1.png}{92.3}{0}
\qfigdef{p1_phy_s1_q04_B}{TOAS - Test 06 (17 Apr 2022)-img/p1_phy_s1_q04_B.png}{128.5}{0}
\qfigdef{p1_phy_s2_q06_1}{TOAS - Test 06 (17 Apr 2022)-img/p1_phy_s2_q06_1.png}{70.8}{0}
\qfigdef{p1_phy_s2_q07_1}{TOAS - Test 06 (17 Apr 2022)-img/p1_phy_s2_q07_1.png}{176.8}{0}
\qfigdef{p1_phy_s2_q08_1}{TOAS - Test 06 (17 Apr 2022)-img/p1_phy_s2_q08_1.png}{181.5}{0}
\qfigdef{p1_phy_s2_q09_1}{TOAS - Test 06 (17 Apr 2022)-img/p1_phy_s2_q09_1.png}{212.0}{0}
\qfigdef{p1_phy_s2_q10_1}{TOAS - Test 06 (17 Apr 2022)-img/p1_phy_s2_q10_1.png}{132.5}{0}
\qfigdef{p1_phy_s2_q11_1}{TOAS - Test 06 (17 Apr 2022)-img/p1_phy_s2_q11_1.png}{81.3}{0}
\qfigdef{p1_phy_s2_q12_1}{TOAS - Test 06 (17 Apr 2022)-img/p1_phy_s2_q12_1.png}{104.1}{0}
\qfigdef{p1_phy_s3_q03_1}{TOAS - Test 06 (17 Apr 2022)-img/p1_phy_s3_q03_1.png}{273.8}{1}
\qfigdef{p1_phy_s3_q04_1}{TOAS - Test 06 (17 Apr 2022)-img/p1_phy_s3_q04_1.png}{180.5}{0}
\qfigdef{p1_phy_s3_q05_1}{TOAS - Test 06 (17 Apr 2022)-img/p1_phy_s3_q05_1.png}{141.6}{0}
\qfigdef{p1_phy_s3_q06_1}{TOAS - Test 06 (17 Apr 2022)-img/p1_phy_s3_q06_1.png}{113.1}{0}
\qfigdef{p1_chem_s1_q03_1}{TOAS - Test 06 (17 Apr 2022)-img/p1_chem_s1_q03_1.png}{93.0}{0}
\qfigdef{p1_chem_s1_q03_A}{TOAS - Test 06 (17 Apr 2022)-img/p1_chem_s1_q03_A.png}{45.8}{0}
\qfigdef{p1_chem_s1_q03_B}{TOAS - Test 06 (17 Apr 2022)-img/p1_chem_s1_q03_B.png}{47.2}{0}
\qfigdef{p1_chem_s1_q03_C}{TOAS - Test 06 (17 Apr 2022)-img/p1_chem_s1_q03_C.png}{43.8}{0}
\qfigdef{p1_chem_s1_q03_D}{TOAS - Test 06 (17 Apr 2022)-img/p1_chem_s1_q03_D.png}{56.5}{0}
\qfigdef{p1_chem_s2_q07_C}{TOAS - Test 06 (17 Apr 2022)-img/p1_chem_s2_q07_C.png}{114.1}{0}
\qfigdef{p1_chem_s2_q07_D}{TOAS - Test 06 (17 Apr 2022)-img/p1_chem_s2_q07_D.png}{189.6}{0}
\qfigdef{p1_chem_s2_q08_1}{TOAS - Test 06 (17 Apr 2022)-img/p1_chem_s2_q08_1.png}{69.2}{0}
\qfigdef{p1_chem_s2_q08_2}{TOAS - Test 06 (17 Apr 2022)-img/p1_chem_s2_q08_2.png}{163.0}{0}
\qfigdef{p1_chem_s3_q03_1}{TOAS - Test 06 (17 Apr 2022)-img/p1_chem_s3_q03_1.png}{199.9}{0}
\qfigdef{p1_chem_s3_q04_1}{TOAS - Test 06 (17 Apr 2022)-img/p1_chem_s3_q04_1.png}{147.4}{0}
\qfigdef{p1_chem_s3_q05_1}{TOAS - Test 06 (17 Apr 2022)-img/p1_chem_s3_q05_1.png}{143.7}{0}
\raggedcolumns
\makeatletter\@beginparpenalty=10000\makeatother
\newcolumntype{L}[1]{>{\raggedright\arraybackslash}p{#1}}
\begin{document}
\begin{center}{\Large\bfseries ALLEN TOAS --- Test 06 (17 Apr 2022)}\end{center}
\begin{multicols}{2}
\raggedcolumns
\partsec{1}{PHYSICS}
\sect{SECTION-I(i)}
\qn{1} A solid non-conducting hemisphere of radius $R$ has a uniformly distributed positive charge of volume density $\rho$. A negatively charged particle having charge $q$ is transferred from centre of its base to infinity. Work performed in the process is :
\begin{qoptions}
\opt{A}{$\dfrac{q\rho R^{2}}{4\in_0}$}
\opt{B}{$\dfrac{q\rho R^{2}}{2\in_0}$}
\opt{C}{$\dfrac{q\rho R^{2}}{8\in_0}$}
\opt{D}{$\dfrac{-q\rho R^{2}}{8\in_0}$}
\end{qoptions}
\qn{2} The potential of electrostatic field is given by $\phi=xy$, where $x$ \& $y$ are coordinate. Which of the following is correct ?
\begin{qoptions}
\opt{A}{Electric field at all points on positive $x$-axis is in positive $y$-direction.}
\opt{B}{Electric field at all points of positive $y$-axis is in negative $x$-direction.}
\opt{C}{Electric field at $(2,2)$ is $\sqrt{2}\,$N/m.}
\opt{D}{Electric field at $x,y$ makes an angle $\theta=\tan^{-1}\!\left(\dfrac{y}{x}\right)$ with $x$-axis.}
\end{qoptions}
\qn{3} Consider a system of $n$ concentric conducting shells of radii $r,2r,4r,8r,\ldots..,2^{n-1}r$ respectively. If inner shell is earthed and $Q$ charge is given on outer most shell, charge appears on inner most shell is :-
\qfig{p1_phy_s1_q03_1}
\begin{qoptions}
\opt{A}{$-Q$}
\opt{B}{$-\dfrac{Q}{2^{n}}$}
\opt{C}{$-\dfrac{Q}{2^{n-1}}$}
\opt{D}{$+Q$}
\end{qoptions}
\qn{4} The bar has a length $\ell$ and mass $m$. It is supported at its ends by rollers of negligible mass. If it is given a small displacement and released, determine the natural frequency of vibration.
\qfig{p1_phy_s1_q04_B}
\begin{qoptions}
\opt{A}{$\omega_N=\sqrt{\dfrac{3g\left(4R^{2}-\ell^{2}\right)}{6R^{2}-\ell^{2}}}$}
\opt{B}{$\omega_N=\sqrt{\dfrac{3g\left(4R^{2}-\ell^{2}\right)^{\frac12}}{6R^{2}+\ell^{2}}}$}
\opt{C}{$\omega_N=\sqrt{\dfrac{3g\left(4R^{2}-\ell^{2}\right)^{\frac12}}{6R^{2}-\ell^{2}}}$}
\opt{D}{$\omega_N=\sqrt{\dfrac{3g\left(4R^{2}+\ell^{2}\right)^{\frac12}}{6R^{2}-\ell^{2}}}$}
\end{qoptions}
\sect{SECTION-I(ii)}
\qn{5} Which of the following relations for a coaxial lens-system separating two media of refractive indices $\mu_1$ and $\mu_2$ are true ? (Here $u$ and $v$ are the distances of the object and the image from the first and the second principal points respectively, $f_1$ and $f_2$ the focal lengths of the system and $m$ is linear transverse magnification.)
\begin{qoptions}
\opt{A}{$\dfrac{f_1}{u}+\dfrac{f_2}{v}=1$}
\opt{B}{$m=-\dfrac{v}{u}\left(\dfrac{f_1}{f_2}\right)$}
\opt{C}{$\dfrac{f_1}{f_2}=-\dfrac{\mu_1}{\mu_2}$}
\opt{D}{$\dfrac{f_1}{u}-\dfrac{f_2}{v}=1$}
\end{qoptions}
\qn{6} Initially two plates (A and B) are given charges as shown. Find the final charges on the plates after the key $k$ is closed. All plates are conducting, parallel and of infinite length and breadth.
\qfig{p1_phy_s2_q06_1}
\begin{qoptions}
\opt{A}{Final charges on plate A is $Q$}
\opt{B}{Final charge on plate B is $0$}
\opt{C}{Final charge on plate C is $3Q/2$}
\opt{D}{Final charge on plate D is $0$}
\end{qoptions}
\qn{7} The 0.8-Mg car travels over the hill having the shape of a parabola. When the car is at point A, it is traveling at 9 m/s and increasing its speed at 3 m/s$^2$. Choose the correct option(s): (Neglect the size of the car)
\qfig{p1_phy_s2_q07_1}
\begin{qoptions}
\opt{A}{The resultant normal force is 6.86 kN}
\opt{B}{The resultant normal force is 2.73 kN}
\opt{C}{Resultant frictional force that all the wheels of the car exert on the road at this instant is 2.22 kN.}
\opt{D}{Resultant frictional force that all the wheels of the car exert on the road at this instant is 1.18 kN.}
\end{qoptions}
\qn{8} The 8-kg sack slides down the smooth ramp. If it has a speed of 1.5 m/s when $y=0.2$ m, choose the correct option(s) :
\qfig{p1_phy_s2_q08_1}
\begin{qoptions}
\opt{A}{Normal reaction the ramp exerts on the sack is 78.44 N}
\opt{B}{Normal reaction the ramp exerts on the sack is 81.84 N}
\opt{C}{The rate of increase in the speed of sack at this instant is 1.96 m/s$^2$}
\opt{D}{The rate of increase in the speed of sack at this instant is 3.92 m/s$^2$}
\end{qoptions}
\qn{9} A tube with reflecting inner walls has a point light source at the bottom of it (see figure). The inner diameter of the tube is $d=12$ mm, the tube's length is $\ell=60$ mm. A convex lens with a focal length $F=36$ mm is placed against the open end of the tube and at a distance $L=90$ mm from the tip of the tube there is a screen, which has a graph paper attached to it. The intersection point with the optical axis O is marked on the paper.
\qfig{p1_phy_s2_q09_1}
\begin{qoptions}
\opt{A}{A point image is formed on the screen at point O.}
\opt{B}{The pattern of images obtained on the screen are parallel straight lines.}
\opt{C}{The pattern of images obtained on the screen are concentric circles.}
\opt{D}{The radius of concentric circles formed on the screen are $(18n)$ mm, where $n$ is an integer.}
\end{qoptions}
\qn{10} Two identical discs, each of negligible mass and radius $r(\lll\ell)$ are connected by a uniform rod of mass $m$ and length $\ell$ as shown in the figure. The disc A is connected with an ideal spring of spring constant $k$ at its centre. The spring is unstretched when $\theta=0^\circ$. There is no slipping at any contact at any point of time during the motion. Choose the correct statement(s) [Spring always obeys Hook's Law and $mg=k\,\ell$]. During the motion, each disc remains in contact with the contact surface :-
\qfig{p1_phy_s2_q10_1}
\begin{qoptions}
\opt{A}{The value of $\theta$, when the system is in stable equilibrium is $60^\circ$}
\opt{B}{The value of $\theta$, when the system is in stable equilibrium is $30^\circ$}
\opt{C}{Spring potential energy of the system in the state of stable equilibrium is $\dfrac{k\ell^{2}}{8}$}
\opt{D}{The angular speed of the rod when it passes through the position of the stable equilibrium is $\sqrt{\dfrac{3g}{4\ell}}$}
\end{qoptions}
\qn{11} A dipole of dipole moment $\vec{P}$ is placed in front of a point charge as shown. Choose the correct statement(s) :
\qfig{p1_phy_s2_q11_1}
\begin{qoptions}
\opt{A}{The dipole will experience a net force irrespective of value of $\theta$.}
\opt{B}{The dipole will experience a net torque about its center irrespective of value of $\theta$.}
\opt{C}{Force experience by dipole is given by $\dfrac{kQP}{r^{3}}\sqrt{1+3\cos^{2}\theta}\,$.}
\opt{D}{If charge $Q$ is not fixed, the whole system can't be in equilibrium for any value of $r$ \& $\theta$.}
\end{qoptions}
\qn{12} Diagram shows a semi-infinite wire having a charge density $\lambda$. A ring carrying a total charge $Q$ is kept as shown. Choose the correct option(s) :
\qfig{p1_phy_s2_q12_1}
\begin{qoptions}
\opt{A}{The force experienced by ring is in vertical direction \& have a value of $\dfrac{\lambda Q}{4\pi\in_0R}$}
\opt{B}{Net force experienced by ring is in vertical direction and have value of $\dfrac{\lambda Q}{2\pi\in_0R}$}
\opt{C}{Tension in the ring is zero.}
\opt{D}{Tension in the ring is $\dfrac{Q\lambda}{8\pi^{2}\in_0R}$}
\end{qoptions}
\sect{SECTION-II}
\qn{1} Two point light sources are positioned on different points of a convex lens' optical axis. The images of these light sources formed by the lens coincide. It is known that one light source is at a distance $a=18$ cm from the center of the lens. How far (in cm) is the other light source distanced from this light source? The focal length of the lens is $f=9$ cm.
\qn{2} Find the coulomb attraction (in N) on a charge of magnitude $5\times10^{-4}$C at the origin due to a charge (of constant density $\rho$) occupying the volume inside the cone $z^2=x^2+y^2$, $0<z<h$. (Take $\rho=\dfrac{\sqrt2}{\pi}\mu$C/m$^3$ , $h=(\sqrt2+1)$ m)
\qfullbegin
\qn{3} Design considerations for the bumper B on the 5000 kg train car require use of a nonlinear spring having the load-deflection characteristics shown in the graph. Select the proper value of $k$ (in MN/m$^2$) so that the maximum deflection of the spring is limited to 0.2 m when the car, traveling at 4 m/s, strikes the rigid stop. Neglect the mass of the car wheels.
\qfig{p1_phy_s3_q03_1}
\qfullend
\qn{4} A string is wrapped over a fixed horizontal cylinder of radius $\dfrac{7}{66}$ m and its other end is connected to a point mass, which hangs in air such that 1.2 m length of string is free. Find minimum velocity (in m/s) given to mass perpendicular to length of cylinder, so that string never gets slack before wrapping over cylinder completely once. (Take $g=9.8$ m/s$^2$)
\qfig{p1_phy_s3_q04_1}
\qn{5} A point object is located at a distance of 100 cm from a screen. A lens of focal length 23 cm mounted on a movable frictionless stand is kept between the source and the screen. The stand is attached to a spring of natural length 50 cm and spring constant 800 N/m as shown. Mass of the stand with lens is 2 kg. How much minimum impulse $P$ (in kg m/s) should be imparted to the stand so that a real image of the object is formed on the screen four times during a complete oscillation. All these are separated by same time interval. (Neglect width of the stand.)
\qfig{p1_phy_s3_q05_1}
\qn{6} Electric field in a region is given by $\vec E=3y\hat i+3x\hat j+5\hat k$. Find out total charge enclosed in a cube of side `10cm' and sides along $x$, $y$ \& $z$ axis as shown :-
\qfig{p1_phy_s3_q06_1}
\raggedcolumns
\partsec{2}{CHEMISTRY}
\sect{SECTION-I(i)}
\qn{1} In a hypothetical reaction, $\mathrm{A(aq)\rightleftharpoons 2B(aq)+C(aq)}$ ($1^{\mathrm{st}}$ order decomposition). `A' is optically active (dextro-rototory) while `B' and `C' are optically inactive but `B' takes part in a titration reaction (fast reaction) with $\mathrm{H_2O_2}$. Hence, the progress of reaction can be monitored by measuring rotation of plane polarised light or by measuring volume of $\mathrm{H_2O_2}$ consumed in titration. In an experiment the optical rotation was found to be $\theta=40^\circ$ at $t=20$ min and $\theta=10^\circ$ at $t=50$ min. from start of the reaction. If the progress would have been monitored by titration method, volume of $\mathrm{H_2O_2}$ consumed at $t=15$ min. (from start) is 40 ml then volume of $\mathrm{H_2O_2}$ consumed at $t=60$ min will be :
\begin{qoptions}
\opt{A}{60 ml}
\opt{B}{75 ml}
\opt{C}{52.5 ml}
\opt{D}{90 ml}
\end{qoptions}
\qn{2} Virial equation of state can be written in terms of molar volume ($V_m$) or in terms of pressure as
\[ Z=1+\frac{B}{V_m}+\frac{C}{V_m^2}+\frac{D}{V_m^3}+\ldots\ldots\ (1) \]
or \quad $Z=1+B'p+C'P^2+\ldots\ldots$ (2)\\
Find $C'$ in terms of $B$ and $C$ at temperature $T$.
\begin{qoptions}
\opt{A}{$C'=\dfrac{C-B^2}{(RT)^2}$}
\opt{B}{$C'=\dfrac{C-B}{(RT)^2}$}
\opt{C}{$C'=\dfrac{C}{RT}$}
\opt{D}{$C'=\dfrac{C}{(RT)^2}$}
\end{qoptions}
\qn{3} Identify the major product :
\qfig{p1_chem_s1_q03_1}
\begin{qoptions}
\opt{A}{\qfigi{p1_chem_s1_q03_A}}
\opt{B}{\qfigi{p1_chem_s1_q03_B}}
\opt{C}{\qfigi{p1_chem_s1_q03_C}}
\opt{D}{\qfigi{p1_chem_s1_q03_D}}
\end{qoptions}
\qn{4} Which of the following complexes has same number of stereoisomeric structures with any geometry (tetrahedral or square planar) ?
\begin{qoptions}
\opt{A}{$\mathrm{[M(NH_3)_2Cl_2]}$}
\opt{B}{$\mathrm{[M(NH_3)(py)(CN)(Cl)]}$}
\opt{C}{$\mathrm{[M(gly)_2]}$}
\opt{D}{$\mathrm{[M(H_2O)_2Br_2]}$}
\end{qoptions}
\sect{SECTION-I(ii)}
\qn{5} One gram of commercial $\mathrm{AgNO_3}$ sample is dissolved in 50 ml of water. If is treated with 50 ml of a KI solution. The silver iodide thus precipitated is filtered off. Excess of KI in the filtrate is titrated with (M/10) $\mathrm{KIO_3}$ solution in presence of 6 M HCl till all $\mathrm{I^-}$ ions are converted into ICl. It requires 50 mL of (M/10) $\mathrm{KIO_3}$ solution, 20 mL of the same stock solution of KI requires 30 mL of (M/10) $\mathrm{KIO_3}$ under similar conditions.
\begin{qoptions}
\opt{A}{The \% of $\mathrm{AgNO_3}$ in the sample $=50\%$}
\opt{B}{The \% of $\mathrm{AgNO_3}$ in the sample $=85\%$}
\opt{C}{Mass of $\mathrm{AgNO_3}$ in the sample $=0.5$ g}
\opt{D}{Mass of $\mathrm{AgNO_3}$ in the sample $=0.85$ g}
\end{qoptions}
\qn{6} In an FCC unit cell a cube is formed by joining the centers of all the tetrahedral voids to generates a new cube. Then the new cube would contain voids as :
\begin{qoptions}
\opt{A}{1 full tetrahedral void}
\opt{B}{1 full octahedral void}
\opt{C}{8 full tetrahedral void and 1 full octahedral void}
\opt{D}{1 full octahedral void only}
\end{qoptions}
\qn{7} Which of following is/are correct :
\begin{qoptions}
\opt{A}{$\mathrm{CO_2}$, $\mathrm{Na^+}$, $\mathrm{AlCl_3}$ all are electrophiles.}
\opt{B}{$\mathrm{F^{\ominus}<Cl^{\ominus}<Br^{\ominus}<I^{\ominus}}$ : nucleophillicity order (in aqueous medium)}
\opt{C}{\qfigi{p1_chem_s2_q07_C}}
\opt{D}{\qfigi{p1_chem_s2_q07_D}}
\end{qoptions}
\qn{8} Correct statement(s) about this compound :
\qfig{p1_chem_s2_q08_1}
\begin{qoptions}
\opt{A}{It shows geometrical isomerism}
\opt{B}{It has centre of symmetry}
\opt{C}{It has plane of symmetry}
\opt{D}{It shows optical isomerism}
\end{qoptions}
\qn{9} \qfigi{p1_chem_s2_q08_2}
\begin{qoptions}
\opt{A}{Total 8 monobrominated product are form}
\opt{B}{Total 5 monobrominated product are form excluding stereoisomer}
\opt{C}{6 monobrominated are optically active}
\opt{D}{4 enantiomeric pair are form}
\end{qoptions}
\qn{10} Which of the following statement(s) is/are \textbf{INCORRECT}?
\begin{qoptions}
\opt{A}{$\mathrm{[Pt(PPh_3)_4]}$ is diamagnetic}
\opt{B}{$\mathrm{[Pt(PPh_3)_4]}$ is $dsp^2$ hybridised}
\opt{C}{$\mathrm{XeO_4}$ is a square planar}
\opt{D}{$\mathrm{XeO_4}$ and $\mathrm{[Pt(PPh_3)_4]}$ have same structure}
\end{qoptions}
\qn{11} Which of the following minerals do/does not contain iron ?
\begin{qoptions}
\opt{A}{Rutile}
\opt{B}{Ilmenite}
\opt{C}{Cerrusite}
\opt{D}{Pyrolusite}
\end{qoptions}
\qn{12} Compound $\mathrm{[Co(gly)_2(CN)(Br)]^+}$ ion has :
\begin{qoptions}
\opt{A}{only one optically inactive isomer}
\opt{B}{two linkage isomers}
\opt{C}{six geometrical isomers}
\opt{D}{eleven stereoisomers}
\end{qoptions}
\sect{SECTION-II}
\qn{1} Gaseous nitrosyl chloride and nitrogen are taken in a flask sealed and heated to some temperature where the total pressure would have been 1 bar if the following equilibrium had not been established.\\
$\mathrm{2NOCl(g)\overset{K_p}{\rightleftharpoons}2NO(g)+Cl_2(g)}$\\
But the actual pressure was found to be 1.2 bar due to above equilibrium .Now into the above equilibrium mixture, some $\mathrm{Cl_2}$ gas was introduced at constant $V$ \& $T$ so that the total pressure would have been 8.3 bar if no further reaction had occurred but the actual pressure was found to be 8.2 bar. find the value of $k_p$.
\qn{2} At 300K, 40 ml of $\mathrm{O_3}$(g) dissolves in 100 g of water at 1.0 atm. What mass of ozone (in gram) dissolved in 1600 g of water at a pressure of 4.0 atm at 300 K ?\\
\relax[Report your answer nearest integer value.]
\qn{3} \qfigi{p1_chem_s3_q03_1}\\
Find number of fraction obtained after fractional distillation of mono-brominated products
\qn{4} \qfigi{p1_chem_s3_q04_1}\\
Degree of unsaturation in P
\qn{5} \qfig{p1_chem_s3_q05_1}
if\\
the Crystal Field Stabilisation Energy of X is $=a*\Delta_o$\\
the Crystal Field Stabilisation Energy of Y is $=b*\Delta_o$\\
the Crystal Field Stabilisation Energy of Z is $=c*\Delta_o$\\
(Neglect pairing energy)\\
then the value of $(a-b)/c$ is
\qn{6} A sample of ore contains\\
(a) iron oxide\\
(b) hydrated alumina\\
(c) $\mathrm{SiO_2}$\\
(d) $\mathrm{TiO_2}$\\
(e) CuO\\
(f) MgO\\
(g) NiO\\
(h) ZnO\\
Number of components soluble in excess $\mathrm{NaOH_{aq.}}$
\raggedcolumns
\partsec{3}{MATHEMATICS}
\sect{SECTION-I(i)}
\qn{1} If $f:[-2\pi,2\pi]\to\mathbb R$ is given by $f(x)=\dfrac{x\cos x}{\left[\frac{x}{5\pi}\right]+\frac12}$ where $[x]$ denotes greatest integer function then $f(x)$
\begin{qoptions}
\opt{A}{is an even function}
\opt{B}{is an odd function}
\opt{C}{is neither even nor odd}
\opt{D}{satisfies $f(x)-f(-x)=2$}
\end{qoptions}
\qn{2} The value of $\displaystyle\lim_{x\to1}\left(\frac{2018}{1-x^{2018}}-\frac{2019}{1-x^{2019}}\right)$ is
\begin{qoptions}
\opt{A}{0}
\opt{B}{$-\dfrac12$}
\opt{C}{$\dfrac12$}
\opt{D}{does not exist}
\end{qoptions}
\qn{3} A circle is drawn circumscribing the square with vertices A, B, C, D (in order) represented by $|x+y|+|x-y|=2$. If tangents at A and B to this circle meets at P, then $\mathrm{PA}^2+\mathrm{PC}^2$ is equal to
\begin{qoptions}
\opt{A}{4}
\opt{B}{8}
\opt{C}{12}
\opt{D}{20}
\end{qoptions}
\qn{4} In triangle ABC, vertex A is $(1,1)$ and internal angle bisector $x-y=2$ of angle B meets the perpendicular bisector of side AC at $(4,2)$. If equation of side AC is $2x+y=3$ and slope of side BC is $m$, then $|4m|$ is equal to
\begin{qoptions}
\opt{A}{4}
\opt{B}{3}
\opt{C}{6}
\opt{D}{7}
\end{qoptions}
\sect{SECTION-I(ii)}
\qn{5} If the number of solutions of equations $\sin^{-1}x+|x|=1$, $\tan^{-1}x+|x|=1$ and $\sec^{-1}x+|x|=1$ are $n_1$, $n_2$ and $n_3$ respectively, then
\begin{qoptions}
\opt{A}{$n_1=1$}
\opt{B}{$n_2=1$}
\opt{C}{$n_2=2$}
\opt{D}{$n_3=1$}
\end{qoptions}
\qn{6} A circle is tangent to line $3x-4y+1=0$ at $(1,1)$ and is normal to line $2x+y-5=0$. Choose correct option(s)?
\begin{qoptions}
\opt{A}{Length of tangent to circle from $(-1,1)$ is 4 units}
\opt{B}{Length of perpendicular from centre of circle to $x$-axis is 3 units.}
\opt{C}{$x$-intercept of director circle of given circle is $\sqrt{82}$}
\opt{D}{$x$-intercept of director circle of given circle is $\sqrt{41}$}
\end{qoptions}
\qn{7} Pairs of lines $\sqrt3(x-2)^2-4(x-2)(y-2)+\sqrt3(y-2)^2=0$ \quad and \quad $\sqrt3x^2-Kxy+\sqrt3y^2=0$; $K\in[2\sqrt3,\,4]$ intersect to form a quadrilateral of area $A_K$.
\begin{qoptions}
\opt{A}{Range of $A_K\in[1,\ \sqrt3]$}
\opt{B}{Quadrilateral is rhombus $\forall\ K\in[2\sqrt3,\ 4]$}
\opt{C}{Range of $A_K\in[0,\ 8-4\sqrt3]$}
\opt{D}{Quadrilateral is kite $\forall\ K\in[2\sqrt3,\ 4]$}
\end{qoptions}
\qn{8} Let $S=\{(x,y)\mid x^2+y^2\geqslant4\}$ and $T=\{(x,y)\mid x^2-y^2\leqslant0\}$. Area of the region given by the set $T-S$ (in sq. units) is
\begin{qoptions}
\opt{A}{$\pi$}
\opt{B}{$2\pi$}
\opt{C}{$3\pi$}
\opt{D}{$4\pi$}
\end{qoptions}
\qn{9} Let two sides AB and BC of a $\Delta$ABC be represented by $x-7y+3=0$ and $x+y+3=0$ respectively. Incentre I $(-3-3\alpha,\ \alpha)$, $\alpha\in R$ of the triangle lies on the altitude $x-y-1=0$ through A then
\begin{qoptions}
\opt{A}{The equation of a line representing side AC is $7x-y=11$}
\opt{B}{The equation of a line representing side AC is $7x+y=1$}
\opt{C}{The area of $\Delta$ABC is $\dfrac{32}{3}$ sq. unit}
\opt{D}{The area of $\Delta$ABC is $\dfrac{32}{5}$ sq.unit}
\end{qoptions}
\qn{10} Let $f(x)=x^{\frac{1}{\{x\}}}$ (where $\{x\}$ denotes fractional part of $x$) then
\begin{qoptions}
\opt{A}{$\displaystyle\lim_{x\to0^+}f(2013^x)=2013$}
\opt{B}{$\displaystyle\lim_{x\to0^-}f(2013^x)=1$}
\opt{C}{$\displaystyle\lim_{x\to0^+}f(2013^x)=e$}
\opt{D}{$\displaystyle\lim_{x\to0}f(2013^x)$ does not exist}
\end{qoptions}
\qn{11} Which of the following statements is/are always true ?
\begin{qoptions}
\opt{A}{number of roots of equation $\displaystyle\sum_{r=1}^{2017}\frac{r}{x-r}=0$ is 2016}
\opt{B}{number of roots of equation $\displaystyle\sum_{r=1}^{2017}\frac{r}{x-r}=0$ is 2017}
\opt{C}{if $f(x)$ is continuous in $[0,2]$ such that $f(0)=f(2)$, then there exists atleast one $x\in[0,1]$ such that $f(x)=f(x+1)$.}
\opt{D}{if $f(x)$ is continuous in $[0,2]$ such that $f(0)=f(2)$, then there exists atleast one $x\in(0,1)$ such that $f(x)=f(x+1)$.}
\end{qoptions}
\qn{12} If a derivable function $y=f(x)$ satisfies the condition $f(x)+f(y)=f(x)f(y)+f(xy)\ \forall x,y\in R$ where $f(1)=0$ and $f'(1)=-2$, then the number of solution(s) of the equation $f(x)=2$ is less than or equal to
\begin{qoptions}
\opt{A}{4}
\opt{B}{3}
\opt{C}{2}
\opt{D}{1}
\end{qoptions}
\sect{SECTION-II}
\qn{1} Number of points where \[ f(x)=\left|(2^x-1)(3^x-3)\,x(x-1)(x-2)\right| \] is not differentiable, is
\qn{2} Let $P(x)=x^{10}+ax^8+bx^6+cx^4+dx^2-1$ be polynomial with real coefficients. If $P(1)=1$, $P(2)=-1$ then the minimum number of real zeroes of $P(x)$ is
\qn{3} If the sum of $n$ terms of the series $\tan^{-1}3+\tan^{-1}\dfrac92+\tan^{-1}\dfrac{33}{4}+\tan^{-1}\dfrac{129}{8}+\ldots.$ is $(2n+1)k\pi-\tan^{-1}2^n$, then $k$ is
\qn{4} $f(x)$ is polynomial which satisfies $|f(x)|\le|e^x-1|\ \forall\ x\in R$ and $A=\{x:x$ is all possible values of $f'(0)\}$, then number of integer in set $A$ is
\qn{5} Given two circle $x^2+y^2+3\sqrt2(x+y)=0$ and $x^2+y^2+5\sqrt2(x+y)=0$, the radius of third circle which is tangent to the two given circles and to their common diametre is
\qn{6} If $f:R-\{0,1\}\to R$ be a function satisfying $f(x)+f\left(1-\dfrac1x\right)=\tan^{-1}x$ and $N=f\left(\dfrac12\right)$, then the value of $\left(\dfrac{8N}{\pi}\right)$ is
\end{multicols}
\end{document}
